Un club d'escalade propose à ses 100 adhérents deux séances par semaine : lundi, jeudi. À chacune des séances, chaque adhérent est libre de venir ou pas.

Le tableau ci-dessous récapitule les choix des adhérents une semaine donnée.

\begin{center}
\begin{tabular}{|c|c|c|c|}
\hline
 & Présent le JEUDI & Absent le JEUDI & Total \\
\hline
Présent le LUNDI & 45 & $x$ & 75 \\
\hline
Absent le LUNDI & 20 & 5 & 25 \\
\hline
Total & 65 & 35 & 100 \\
\hline
\end{tabular}
\end{center}

\underline{Exemple :} le tableau montre que 45 adhérents sont venus lundi et jeudi.

\begin{enumerate}
\item Décrire par une phrase ce que représente le nombre $x$ et déterminer sa valeur.

\item On choisit un adhérent au hasard.
\begin{enumerate}
\item Quelle est la probabilité qu'il s'agisse d'un adhérent qui n'est venu ni le lundi ni le jeudi ?
\item Quelle est la probabilité qu'il s'agisse d'un adhérent qui n'est venu qu'un seul jour ?
\item On sait à présent que l'adhérent choisi est venu le lundi. \\
Quelle est la probabilité qu'il soit également venu le jeudi ?
\end{enumerate}

\item Chacun des adhérents verse au club une cotisation annuelle de 100 euros.
\begin{enumerate}
\item En 2026, le club compte 100 adhérents. \\
Quel est le montant total des cotisations versées au club en 2026 ?
\item On suppose que, de 2026 (inclus) à 2041 (inclus) le montant de la cotisation reste stable, mais que le nombre d'adhérents augmente régulièrement de 5 unités chaque année. Ainsi, en 2026, il y a 100 adhérents, en 2027, il y a 105 adhérents, en 2028, il y a 110 adhérents, en 2029, il y a 115 adhérents, etc.

Quel sera le montant total des cotisations versées au club entre 2026 et 2041 ?

\underline{Indication :} on pourra utiliser la formule ci-dessous :

\[a + (a + r) + (a + 2r) + (a + 3r) + \cdots + (a + nr) = \frac{2a + nr}{2} \times (n + 1).\]
\end{enumerate}
\end{enumerate}
